\documentclass[10pt,a4paper]{amsart}
\usepackage[margin=19mm]{geometry}
\usepackage{amsmath,amssymb,mathtools}
\usepackage[scr=rsfs,cal=euler]{mathalfa}
\usepackage{xfrac}
\usepackage[unicode,hidelinks,bookmarks=false]{hyperref}
\newcommand{\ie}{i.e.\ }
\newcommand{\cf}{cf.\ }
\newcommand{\resp}{respectively\ }
\newcommand{\Subsection}{\subsection}
\providecommand{\nobreakdash}{\nobreak\hbox{-}}
\setlength{\emergencystretch}{2em}
\multlinegap0pt
\usepackage[shortlabels]{enumitem}
\usepackage{mathtools}
\usepackage{amssymb}
\usepackage{xcolor}
\usepackage{algorithm}\usepackage{algpseudocode}
\usepackage{tikz}
\newtheorem{theorem}{Theorem}
\title{Chemotactic Feedback Controls Patterning in Hybrid Tumor–Stroma Model}
\author{Jiguang Yu, Louis Shuo Wang, Zonghao Liu, Jingfeng Liu}
\date{Source: arXiv:2601.16337v2 (CC BY 4.0)}
\begin{document}
\maketitle

\noindent\emph{Source and licence.} Chemotactic Feedback Controls Patterning in Hybrid Tumor–Stroma Model. Version arXiv:2601.16337v2 (CC BY 4.0), distributed by arXiv under the Creative Commons Attribution 4.0 International licence, http://creativecommons.org/licenses/by/4.0/, as recorded in the arXiv licence metadata for this exact version and shown on the abstract page https://arxiv.org/abs/2601.16337v2. Full licence notice: \texttt{licenses/arxiv-2601.16337-CC-BY-4.0.txt}.\par\smallskip

\noindent\textit{Editorial scope.} Counted unit: the article's theorem thm:positivity (source0.tex lines 815--832). The article's complete original proof of it is delivered with it (source0.tex lines 834--934, one \texttt{proof} environment of 101 lines). No mathematics is written, repaired or re-derived: the statement and the proof are the article's own lines, in the article's own notation, and no retained line is substituted or re-flowed.  Retained uncounted as the source-local context the proof needs: lines 294--303; lines 549--554.  Nothing was omitted from any retained span, so the package carries no omission record.  No retained line contains a citation, so the wrapper carries no bibliography and no bibliography entry.  No retained line contains a figure environment, an image-inclusion command or drawing code, so no image is delivered and the package declares no asset dependency.  Editorial formatting only, in addition: an AMS \texttt{amsart} preamble that restates the article's own declarations and macros used by the retained lines, namely the environments \texttt{theorem} on separate counters, together with the standard packages those lines need.  The retained environments print in this document as Theorem 1 in order of appearance, whereas the article prints the same items with the numbering of its own full document.  The complete native source of the article as distributed in the arXiv e-print for this exact version is bundled unchanged as \texttt{sources/193228/source0.tex}.
\begin{align}
\label{eq:system}
\begin{cases}
\partial_{t} S = d_{S}\Delta S + \lambda_{S} S \left(1 - \tfrac{S+R}{K}\right) - \alpha S - \delta(I)S + \xi \bigl[1 - \phi(I)\bigr]R, & x\in U,\ t>0, \\[4pt]
\partial_{t} R = d_{R}\Delta R + \lambda_{R} R \left(1 - \tfrac{S+R}{K}\right) + \alpha S + \eta \phi(I)F_aR - \xi \bigl[1 - \phi(I)\bigr]R, & x\in U,\ t>0, \\[4pt]
\partial_{t} I = d_{I}\Delta I - \gamma_{I}I, & x\in U,\ t>0, \\[4pt]
\partial_{t} P = -\theta \phi(I)P + \beta \bigl[1 - \phi(I)\bigr]F_a, & x \in U,\ t>0, \\[4pt]
\partial_{t} F_a = \theta \phi(I)P - \beta \bigl[1 - \phi(I)\bigr]F_a, & x \in U,\ t>0.
\end{cases}
\end{align}

\begin{align}
\label{eq:base_vector_form}
\partial_t\mathbf u-\mathcal D_{\mathrm{diff}}\Delta \mathbf u=\mathcal G(\mathbf u),
\qquad 
\mathcal D_{\mathrm{diff}}=\mathrm{diag}(d_S,d_R,d_I,0,0),
\end{align}

\begin{theorem}[\textbf{Nonnegativity}]
\label{thm:positivity}
Let $U\subseteq\mathbb{R}^N$ ($N\in\{1,2,3\}$) be a bounded domain with $C^\infty$ boundary, and consider the base system \eqref{eq:system} under homogeneous Neumann boundary conditions
\[
\partial_{\mathrm n}S=\partial_{\mathrm n}R=\partial_{\mathrm n}I=0 \qquad \text{on }\partial U\times(0,\infty).
\]
Assume that all parameters in \eqref{eq:system} are nonnegative constants and that the initial data satisfy
\[
S_0,R_0,I_0,P_0,F_{a,0}\ge 0 \quad \text{a.e.\ in }U,
\qquad 
(S_0,R_0,I_0,P_0,F_{a,0})\in (L^\infty(U))^5.
\]
Then any classical solution $\mathbf{u}=(S,R,I,P,F_a)$ of \eqref{eq:system} on its maximal interval of existence $[0,T_{\max})$ remains nonnegative:
\[
S(x,t),\,R(x,t),\,I(x,t),\,P(x,t),\,F_a(x,t)\ge 0
\qquad \text{for all }(x,t)\in \overline U\times[0,T_{\max}).
\]
\end{theorem}

\begin{proof}
\noindent\textbf{Step 1: Quasi-positivity.}
Let $\mathcal G=(G_1,\dots,G_5)^{\mathsf T}$ be as in \eqref{eq:base_vector_form} with $\mathbf u=(S,R,I,P,F_a)^{\mathsf T}$. 
We verify that $\mathcal G$ is quasi-positive on $\mathbb{R}^5_{\ge 0}$: for each component $i$, if $u_i=0$ and $\mathbf u\ge 0$, then $G_i(\mathbf u)\ge 0$. 
Indeed, for $\mathbf u\ge 0$,
\begin{align*}
G_1\big|_{S=0} &= \xi(1-\phi(I))R \ge 0,\\
G_2\big|_{R=0} &= \alpha S \ge 0,\\
G_3\big|_{I=0} &= 0,\\
G_4\big|_{P=0} &= \beta(1-\phi(I))F_a \ge 0,\\
G_5\big|_{F_a=0} &= \theta\phi(I)P \ge 0,
\end{align*}
since $\phi(I)\in(0,1)$ and all parameters are nonnegative.

\smallskip
\noindent\textbf{Step 2: Positivity of the diffusing components $(S,R,I)$.}
Let $\mathbf{u}=(S,R,I)$. We define the maximal interval of nonnegativity as
\begin{equation*}
    \tau\coloneqq \sup\left\{t\in [0,T_{\max}): S(\cdot,s), R(\cdot,s), I(\cdot,s)\ge 0 \text{ on } \bar U \text{ for all } s\in [0,t] \right\}.
\end{equation*}
We claim that $\tau=T_{\max}$. 
Suppose, for the sake of contradiction, that $\tau<T_{\max}$.
By the definition of $\tau$ and continuity, we have:
\begin{enumerate}[label=(\roman*)]
    \item $S(x,t) \ge 0$, $R(x,t)\ge 0$, and $I(x,t)\ge 0$ for all $x\in \bar U$ and $t\in [0,\tau]$.
    \item There exists a component (without loss of generality, assume $S$) and a point $x^*\in \bar U$ such that $S(x^*,\tau)=0$.
\end{enumerate}
At the point $(x^*,\tau)$, $S$ attains a global spatial minimum (value $0$). If $x^*\in U$, then $\Delta S(x^*,\tau)\ge 0$ by the standard necessary condition for a minimum. We next show that the same sign condition holds if $x^*\in \partial U$. The homogeneous Neumann boundary condition,
\[
\partial_{\mathrm n} S(x^*,\tau) = \nabla S(x^*,\tau)\cdot \mathrm n(x^*) = 0,
\]
implies that the gradient of $S$ at $x^*$ lies in the tangent space $T_{x^*}\partial U$, where $\mathrm n(x^*)$ is the unit outward normal vector of $\partial U$ at $x^*$.
Taylor expansion about $S$ at $x^*$ gives
\begin{equation*}
    S(x,\tau) = S(x^*,\tau) + \nabla S(x^*,\tau)\cdot (x-x^*) + \frac{1}{2} (x-x^*)^\top \mathrm{H}S(x^*,\tau) (x-x^*) + o(\lVert x-x^*\rVert ^2), 
\end{equation*}
where $\mathrm{H}S(x^*,\tau)$ is the Hessian matrix of $S$ at $x^*$.
The minimum at $x^*$ implies that either 
\begin{equation}\label{eq:minimum_condition1}
\nabla S(x^*,\tau)\cdot (x-x^*) \ge 0,
\end{equation}
or, 
\begin{equation}\label{eq:minimum_condition2}
\nabla S(x^*,\tau)\cdot (x-x^*) = 0 \text{ and } \frac{1}{2} (x-x^*)^\top \mathrm{H}S(x^*,\tau) (x-x^*)\ge 0
\end{equation}
Let $x\in U$, we define a projection operator:
\begin{align*}
    T_0: U &\to T_{x^*}\partial U \\[4pt]
       x &\mapsto (x-x^*) + \Big(\mathrm n(x^*)\cdot (x-x^*)\Big)\, \mathrm{n}(x^*).
\end{align*}
$T_0(x)$ is the projection of $x-x^*$ onto the tangent space $T_{x^*}\partial U$. Since $U$ is a bounded domain with $C^\infty$ boundary $\partial U$, it is easy to see the linear span of $T_0(x)$ for $x\in U$ is the whole tangent space:
\begin{equation*}
    \operatorname{span}_{\mathbb{R}} \{T_0(x): x\in U\} = T_{x^*}\partial U.
\end{equation*}
\eqref{eq:minimum_condition1} and $\nabla S(x^*,\tau)\in T_{x^*}\partial U$ ensure that $\nabla S(x^*,\tau) = 0$. \eqref{eq:minimum_condition2} ensures that $\mathrm{H}S(x^*,\tau)$ is positive semidefinite. Therefore, we have for $x^*\in \partial U$,
\[
\Delta S(x^*,\tau) \ge 0.
\]
Going back to the $S$-equation at $(x^*,\tau)$:
\[
\partial_t S(x^*,\tau) = d_S\Delta S(x^*,\tau) + G_1(S = 0) \ge 0.
\]
Therefore, $S$ cannot become negative in the immediate future $t>\tau$. This contradiction implies $\tau=T_{\max}$.

\smallskip
\noindent\textbf{Step 3: Positivity of the stromal ODE subsystem $(P,F_a)$.}
For each fixed $x\in U$ the stromal variables satisfy the linear ODE system
\begin{equation}\label{eq:PA_ODE_pointwise}
\frac{d}{dt}
\begin{pmatrix}
P(x,t)\\ F_a(x,t)
\end{pmatrix}
=
B(I(x,t))
\begin{pmatrix}
P(x,t)\\ F_a(x,t)
\end{pmatrix},
\qquad
B(I)=
\begin{pmatrix}
-\theta\phi(I) & \beta(1-\phi(I))\\
\theta\phi(I) & -\beta(1-\phi(I))
\end{pmatrix}.
\end{equation}
Note that $B(I)$ is a Metzler matrix for every $I\ge0$ (its off-diagonal entries are nonnegative).

Fix $x\in U$ and write $p(t)=P(x,t)$, $a(t)=F_a(x,t)$.
If $p(t_0)=0$ at some time $t_0$ with $a(t_0)\ge0$, then from \eqref{eq:PA_ODE_pointwise},
\[
p'(t_0)=\beta(1-\phi(I(x,t_0)))\,a(t_0)\ge 0,
\]
so $p$ cannot decrease below $0$ at $t_0$.
Similarly, if $a(t_0)=0$ with $p(t_0)\ge0$, then
\[
a'(t_0)=\theta\phi(I(x,t_0))\,p(t_0)\ge 0,
\]
so $a$ cannot decrease below $0$ at $t_0$.
Hence the nonnegative quadrant in $(p,a)$ is forward invariant, proving the claim.

Combining Steps 2 and 3 yields nonnegativity of all components on $\overline U\times[0,T_{\max})$.
\end{proof}

\end{document}
